\section{Discussions and summary}
\label{discussion}
\indent \indent We have simulated a three-dimensional model of squaric acid crystal considering the interlayer interaction ($J_3$) of protons that were not rigorously accounted earlier. We use the pseudo-spin formalism and introduce an additional interlayer interaction ($J_3$) of the Ising type. We aim at understanding the role played by the nature of the $J_3$ and the strength of it on the phases that have been found for quasi-2D case. Motivated by earlier works~\cite{matsushita1982cluster}, in the present work, we use three different types of interlayer coupling (AFM, FM, Disorder). Our QMC analysis reveals interesting results for small $J_2$ strengths. Similar to the quasi-2D case~\cite{ishizuka2011quantum,vijigiri2020dipole}, we observe an intermediate-state appearing (around $T\sim 0.3$) where the order parameter, $P$, is found to exhibit a plateau behavior that vanishes with the system size. This is not the scenario for all the types of couplings considered. From Fig.~\ref{ice_param}(a), for an AFM case of interlayer coupling, we see that the AFE stacked configuration of layers seems to be more robust (see Fig.~\ref{order_param}) as the value of $P$($\sim 1$) at low temperatures does not change even when $J_3$ is increased. Since the quantum fluctuations in the AFM case has more pronounced quantum fluctuations, therefore, the intermediate liquid-like states is found to vanish as the strength of $J_3$ is increased (see Fig.~\ref{ice_param}. The plateau starts getting narrowed and then disappears for large $J_3$). Secondly, when we consider the FM type of $J_3$, we observe that not only does the AFE stacked configuration is found to exist for large $J_3$ ($J_3<1.0$), but also the intermediate state was shown to extend to larger temperatures (and fields). In contrast to the AFM coupling of $J_3$, the anomalous peak seen at lower temperatures in the specific-heat curves is shown to have lesser entropy. As the peak height is lesser than in the case of AFM. We know that the entropy ($S\propto \ln(C/T) $). Therefore, the higher the peak higher the more degeneracy the system possesses. Nevertheless, in the last case of disordered coupling, we see neither of the AFE nor intermediate liquid-like state is shown to exist beyond a disorder strength of $J_3=0.4$. This is intuitive because the disorder case accounts for a large degree of quantum fluctuations compared to AFM and FM resulting in such behavior. \\
\indent  We note that previous studies~\cite{vijigiri2020dipole,ishizuka2011quantum} has investigated the finite-temperature phase diagram in a related model with two body interactions~\cite{ishizuka2011quantum} followed by a four-spin model~\cite{vijigiri2020dipole,chern2014gauge} motivated by the earlier studies~\cite{maier1982interacting}. Though it is known that the interlayer interaction is a weak one, some studies have reported that in $\text{H}_2\text{SQ}$ interlayer interaction is more likely AFM in nature. Here, in the present work, we show that AFM type $J_3$ has a more robust anomalous peak compared to the others. And the stability of the intermediate state strictly depends on the strength. While this is the scenario for AFM, the FM coupling also hosts the AFE phase, but with the exception that unlike the AFM coupling the intermediate state extends to larger temperatures and pressures as the strength of $J_3$ is increased. \\
\indent Further, apart from the ground state properties and static susceptibilities, to probe the nature of excitations that can be experimentally verified, we also calculate the dynamic structure factor in our calculations. We use a real-space linear spin-wave theory to obtain the spectrum. We see that the characteristic behavior of the deconfined phase can be seen in Fig.~\ref{density_plot}(a) when $J_{1,2}$ are absent. That is, apart from the flat spectrum in Fig.~\ref{density_plot}(a) the spectrum is highly asymmetric around the $X, M$ high symmetry point. This is expected because in the deconfined phase the system along the two orthogonal directions has an independent configuration. The configuration along the unit vector $x$ is independent of configuration along $y$. Due to this fact, we see that the spectrum until the $M-$high symmetry point is broader while the spectrum from $M$ to $\Gamma$ is flat. This explains the deconfined characteristic. Next, once we put the finite $J_1$ with small $J_2$, we see that a small deviation from the flat spectrum from $M$ to $\Gamma$ is seen. This is because the small $J_2$ is responsible for a global ordering in the system. And in the global ordering case, the ordering along the two orthogonal vectors spanning the lattice is no more independent. This can be further corroborated from the Fig.~\ref{density_plot}(c) when the $J_2$ is further increased a sinusoidal behavior in the spectrum is seen. We know that any ordering in the system manifests in the form of $\sim \cos(\vec{k})$ or $\sin(\vec{k})$. This explains the dynamic structure factor in three important regimes.
